% !TEX root = ../main.tex
%% Rozdział 3. Plan: 6–8 stron.

\chapter{Projekt systemu}
\label{chap:projekt}

\CoTuMaByc{Krótki wstęp: co powstało (biblioteka generatora, okno edytora, aplikacja do badań) i z jakich części się składa.}

\section{Wymagania}
\label{sec:wymagania}

\CoTuMaByc{Wymagania funkcjonalne i niefunkcjonalne wynikające z opisu tematu. Klocek: \kol{tab:wymagania}.}

\section{Architektura}
\label{sec:architektura}

\CoTuMaByc{Cztery biblioteki kodu i zależności między nimi; wspólny rdzeń \kol{WFC.Core} dla okna edytora i aplikacji. Klocki: \kol{tab:biblioteki}, \kol{fig:architektura}.}

\section{Kafelki i reguły}
\label{sec:kafelki}

\CoTuMaByc{Moduł siatki, modele, etykiety ścian, obroty, plik reguł i jego kompilacja. Zestawy o rosnącej złożoności i warianty reguł. Klocki: \kol{tab:parametry-regul}, \kol{tab:zestawy}, \kol{tab:warianty-regul}, \kol{fig:modele}, \kol{fig:orientacja}.}

\section{Warianty generatora}
\label{sec:warianty}

\CoTuMaByc{Wymienne elementy generatora: rozwiązywacze, heurystyki wyboru, reakcje na sprzeczność, sposoby rysowania i ustawienia łączenia rysowania w \ac{urp}. Klocki: \kol{tab:rozwiazywacze}, \kol{tab:heurystyki}, \kol{tab:reakcje}, \kol{tab:rysowanie}, \kol{tab:laczenie}.}

\section{Okno edytora i aplikacja}
\label{sec:aplikacja}

\CoTuMaByc{Editor Tool do projektowania poziomów (podgląd krok po kroku, zapis mapy jako prefab) i aplikacja do badań. Zrzuty ekranu. Klocki: \kol{fig:okno-generatora}, \kol{fig:krok-po-kroku}, \kol{fig:prefab}, \kol{fig:aplikacja}.}
